\chapter{Acromegaly}
\index{Acromegaly}

\section*{Definition and Overview}
Acromegaly is a rare hormonal disorder in which the body produces too much growth hormone, usually because of a small, non-cancerous tumour (adenoma) in the pituitary gland at the base of the brain. The excess hormone causes the bones, soft tissues, and organs to enlarge, particularly the hands, feet, and face.

The changes develop slowly over many years, which is why the condition often goes unrecognised for a long time. When the same excess occurs before the growth plates close in childhood, it causes a different condition, gigantism, in which the child grows unusually tall. The name acromegaly comes from the Greek words for ``extremities'' and ``enlargement'', reflecting the characteristic overgrowth of the hands, feet, and facial bones.

\section*{Epidemiology}
Acromegaly is rare, affecting roughly 1 in 100,000 people at any one time. It is usually diagnosed in adults between the ages of 40 and 60, and men and women are affected about equally.

Because the changes are gradual and non-specific, diagnosis is often delayed by several years. Many people consult doctors for joint pain, sleep problems, headaches, or diabetes before the underlying hormonal cause is recognised, and the diagnosis is sometimes first suspected from a change in photographs or in ring and shoe sizes.

\section*{Aetiology and Pathophysiology}
In more than 9 in 10 cases, acromegaly is caused by a benign tumour of the pituitary gland that produces excess growth hormone. The excess hormone signals the liver to produce more insulin-like growth factor (IGF-1), and this surge drives the enlargement of bones and soft tissues.

Because the tumour sits at the base of the brain, it can press on the optic nerves and cause visual problems, and it can interfere with the pituitary's production of other hormones. Rarely, growth hormone excess comes from a tumour elsewhere in the body, such as in the lung or pancreas. There is no clear preventable cause, and the tumour is usually due to a spontaneous change within a single pituitary cell.

\section*{Clinical Features}
\begin{itemize}
    \item \textbf{Hands and feet:} progressive enlargement, such as the need for bigger rings, gloves, or shoes
    \item \textbf{Facial change:} coarsening of features, with a prominent brow, jaw, and nose, and enlarged lips and tongue
    \item \textbf{Head and neck:} a deepening voice and increased spacing between the teeth
    \item \textbf{Skin changes:} excessive sweating, oily skin, and skin tags
    \item \textbf{Pain and pressure:} headache, joint pain, stiffness, and carpal tunnel syndrome
    \item \textbf{Sleep:} sleep apnoea with loud snoring and daytime sleepiness, plus fatigue and weakness
    \item \textbf{Metabolic and visual:} high blood pressure, diabetes, and visual problems with larger tumours
\end{itemize}

\section*{Differential Diagnosis}
The gradual changes of acromegaly can be confused with other conditions.

\begin{itemize}
    \item Constitutional tall stature or large body habitus without hormonal disease
    \item Gigantism, when growth hormone excess begins before the growth plates close
    \item Hypothyroidism, which can cause facial puffiness and tiredness
    \item Other causes of headaches, joint pain, carpal tunnel syndrome, and sleep apnoea
    \item Paget's disease of bone and other bone disorders
\end{itemize}

\section*{When to Seek Medical Care}
See a doctor if you notice progressive enlargement of your hands, feet, or facial features, a change in ring or shoe size, or a combination of new headaches, excessive sweating, joint pain, and snoring. Any sudden change in vision, or a sudden severe headache with vomiting, needs urgent assessment, because a pituitary tumour can occasionally bleed or swell suddenly.

\textbf{Contact your local emergency services for:}
\begin{itemize}
    \item Sudden, severe headache with vomiting
    \item Sudden loss of vision or double vision
    \item Confusion, weakness, or collapse
    \item Any sudden, severe, or concerning symptom
\end{itemize}

\section*{Diagnosis and Investigations}
\begin{itemize}
    \item \textbf{Blood tests:} measurement of IGF-1, which reflects the average growth hormone level over time
    \item \textbf{Growth hormone suppression test:} growth hormone is measured after a glucose drink; in acromegaly it fails to fall normally
    \item \textbf{MRI of the pituitary:} to find the tumour and assess its size and position
    \item \textbf{Visual field testing:} checks for pressure on the optic nerves
    \item \textbf{Other hormone tests:} assessment of the remaining pituitary hormones
    \item \textbf{Screening for complications:} tests for diabetes, high blood pressure, and, where indicated, heart and bowel problems
\end{itemize}

\section*{Management}
Treatment aims to bring growth hormone and IGF-1 back to normal, control the tumour, and manage its effects on the body.

\begin{itemize}
    \item \textbf{Surgery:} removal of the tumour through the nose (transsphenoidal surgery) is the usual first treatment and can cure many small tumours
    \item \textbf{Medicines:} injections of somatostatin analogues, medicines that block the growth hormone receptor, or dopamine agonists can lower hormone levels
    \item \textbf{Radiotherapy:} used for tumours that persist, recur, or are not controlled by medicines
    \item \textbf{Management of complications:} treatment of high blood pressure, diabetes, sleep apnoea, and joint pain
    \item \textbf{Hormone replacement:} if the pituitary fails after surgery or radiotherapy
    \item \textbf{Long-term monitoring:} regular blood tests, MRI scans, and reviews with an endocrinologist
\end{itemize}

\section*{Complications}
\begin{itemize}
    \item Cardiovascular disease, including high blood pressure, an enlarged heart, and heart failure
    \item Diabetes and impaired glucose tolerance
    \item Obstructive sleep apnoea and chronic fatigue
    \item Arthritis, joint degeneration, and carpal tunnel syndrome
    \item An increased risk of bowel polyps and colorectal cancer, so colonoscopy screening is recommended
    \item Visual loss from pressure on the optic nerves
    \item Hormone deficiencies after surgery or radiotherapy
\end{itemize}

\section*{Prognosis and Outcomes}
With treatment, growth hormone levels can usually be brought under control, and symptoms such as headaches, sweating, and sleep apnoea often improve. Some changes to the bones and face are permanent, but soft-tissue swelling can resolve. Control of hormone levels reduces the risk of serious complications, particularly heart and blood-vessel disease.

The outcome depends on how early the condition is detected and how well hormones are controlled: well-controlled disease allows a near-normal life expectancy, whereas uncontrolled disease shortens it, mainly through cardiovascular complications. Lifelong follow-up with an endocrinologist is therefore recommended.

\section*{Prevention and Self-Care}
\begin{itemize}
    \item Be alert to gradual changes in ring, glove, or shoe size and report them early
    \item Attend all follow-up appointments, including blood tests and MRI scans
    \item Keep blood pressure, blood sugar, and cholesterol under review and follow treatment advice
    \item Discuss sleep problems or snoring with your doctor so that sleep apnoea is not missed
    \item Follow recommended screening for bowel polyps and for heart complications
    \item Maintain a healthy weight, exercise within your abilities, and care for your joints
\end{itemize}

\section*{Sources and Further Reading}
The following organisations provide reliable information on acromegaly for patients and health professionals.

\begin{itemize}
    \item NHS. \textit{Information on acromegaly, diagnosis, and treatment.} https://www.nhs.uk/conditions/
    \item Mayo Clinic. \textit{Guide to acromegaly and its management.} https://www.mayoclinic.org/
    \item MedlinePlus. \textit{Health information on acromegaly and pituitary disorders.} https://medlineplus.gov/
    \item MSD Manual. \textit{Professional information on pituitary and hormonal disorders.} https://www.msdmanuals.com/
    \item NICE. \textit{Clinical guidance relevant to pituitary tumours and endocrine disorders.} https://www.nice.org.uk/
\end{itemize}
